\documentclass[11pt]{article}
\usepackage[a4paper,margin=1in]{geometry}
\usepackage{amsmath,amssymb,mathtools,bm}
\usepackage{enumitem,booktabs,array}
\usepackage{tcolorbox}
\usepackage{hyperref}
\usepackage[protrusion=true,expansion=false]{microtype}
\hypersetup{colorlinks=true,linkcolor=blue,urlcolor=blue,citecolor=blue}
\setlist[itemize]{topsep=2pt,itemsep=2pt}
\setlist[enumerate]{topsep=2pt,itemsep=2pt}
\newtcolorbox{keybox}{colback=blue!3!white,colframe=blue!40!black,boxrule=0.4pt,arc=1mm}
\newtcolorbox{alertbox}{colback=red!3!white,colframe=red!40!black,boxrule=0.4pt,arc=1mm}
\newtcolorbox{answerbox}{colback=green!3!white,colframe=green!40!black,boxrule=0.4pt,arc=1mm}
\newtcolorbox{drillbox}{colback=yellow!5!white,colframe=orange!60!black,boxrule=0.4pt,arc=1mm}
\newcommand{\EV}{\mathbb{E}}
\newcommand{\Var}{\mathrm{Var}}
\newcommand{\Cov}{\mathrm{Cov}}
\newcommand{\blank}[1]{\underline{\hspace{#1}}}

\title{E02-02: Albuquerque, Fos \& Schroth (2022, JFE) \\
\large ``Value creation in shareholder activism'' --- Active Recall Workbook}
\author{Empirical Corporate Finance II --- Exam Prep}
\date{\today}
\begin{document}\maketitle

\begin{keybox}
\textbf{Paper in one sentence.} AFS estimate a structural model of activism decomposing the activist's gain into stock-picking (information advantage) vs.\ value-creation (engagement that raises firm value), and find that the bulk ($>$70--80\%) of the activist's realized return comes from \emph{real value creation}, not from trading on prior information --- validating the monitoring-creates-value view.

\textbf{Placement.} Structural / mechanism companion to Brav-Jiang-Kim (E02-01). Provides the quantitative decomposition that defends activism against the ``tipping / information'' critique.
\end{keybox}

\section*{Round 1: Complete Walk-Through}

\subsection*{1.1 Research question}
When an activist buys a 5\%+ stake and triggers a 13D filing, the target's stock jumps. Is the activist profiting from:
\begin{enumerate}[label=(\roman*)]
\item Pre-existing information advantage (``stock picking''), buying cheap with knowledge of a future catalyst?
\item Engagement that changes firm policy and raises fundamental value (``value creation'')?
\end{enumerate}

\subsection*{1.2 Structural model}
\begin{itemize}
\item Model activist's accumulation period (pre-13D) and post-13D engagement.
\item Latent firm-value process affected by activist engagement with probability $p$.
\item Activist chooses stake $\alpha$ and engagement effort to maximize expected profit; data on prices, volumes, stakes, and outcomes identify $p$, cost of effort, and information parameters.
\end{itemize}

\subsection*{1.3 Data}
\begin{itemize}
\item 13D filings 1993--2014; $\sim$2{,}000 events.
\item CRSP price and volume data around filings.
\item Hand-collected campaign outcomes (board seats, spin-offs, M\&A).
\end{itemize}

\subsection*{1.4 Main quantitative findings}
\begin{itemize}
\item Activists' total profit on successful campaigns averages $\sim$20--30\% of their stake value.
\item Decomposition: $\sim$70--85\% of that profit is \emph{value creation} (engagement improves firm value); $\sim$15--30\% is stock picking.
\item Announcement CARs around 13D filings reflect market's update about value creation.
\item Value-creation share larger for offensive, strategic, and governance campaigns; smaller for financial campaigns.
\end{itemize}

\subsection*{1.5 Mechanisms}
\begin{itemize}
\item Engagement raises operating performance and payout; value-creation channel explicit in the structural model.
\item Stock-picking share exists but is secondary; activists accumulate under constraints (disclosure, liquidity).
\item Complements BJK by attributing BJK's plant-level gains quantitatively to activist engagement.
\end{itemize}

\subsection*{1.6 Policy implications}
\begin{itemize}
\item Activists' private gains are largely socially valuable (real value created).
\item Policy restrictions on disclosure windows should consider trade-off between limiting stock-picking and preserving engagement incentives.
\item Supports rather than refutes the case for activist-investor protections.
\end{itemize}

\subsection*{1.7 Caveats}
\begin{alertbox}
\begin{itemize}
\item Structural estimates depend on functional-form assumptions.
\item Sample is 13D filings; non-public campaigns (``wolf packs'') are not captured directly.
\item Post-2014 activism environment differs (ESG, universal owners) so external validity of point estimates is uncertain.
\end{itemize}
\end{alertbox}

\section*{Round 2: Level 1 Fill-in}
\begin{enumerate}
\item Two channels: stock \blank{1.5cm} vs.\ value \blank{1.5cm}.
\item Data: \blank{1cm}D filings 1993--\blank{1cm}.
\item Activists' profit: $\sim$\blank{1cm}--\blank{1cm}\% of stake.
\item Value-creation share: $\sim$\blank{1cm}--\blank{1cm}\%.
\item Campaigns with largest value-creation: \blank{1.5cm} and governance.
\end{enumerate}

\section*{Round 3: Level 2 Prompts}
\begin{enumerate}
\item Describe the structural decomposition of activist profit.
\item Report the main numerical finding and its interpretation.
\item Relate AFS to Brav-Jiang-Kim (2015).
\item Discuss policy implications for disclosure windows and activism regulation.
\item What assumptions drive the decomposition, and how sensitive are the results?
\end{enumerate}

\section*{Round 4: Minimal Scaffolding}
From a blank page: (i) question, (ii) structural model, (iii) quantitative results, (iv) mechanism, (v) policy.

\section*{Round 5: Blank Page}
Essay: decomposing the returns to shareholder activism --- what AFS adds to BJK.

\vspace{6pt}
\begin{drillbox}
\textbf{Speed Drill.} (1) Two channels. (2) Value-creation share. (3) Data source. (4) Policy implication. (5) Complement to BJK.
\end{drillbox}

\section*{Quick Reference \& Answer Key}
\begin{answerbox}
\begin{itemize}
\item \textbf{Model.} Structural: stock-picking + value-creation.
\item \textbf{Profit share.} $\sim 70$--$85$\% from value creation.
\item \textbf{Stock picking.} $\sim 15$--$30$\%.
\item \textbf{Complement.} AFS quantifies; BJK shows plant-level realization.
\end{itemize}
\end{answerbox}

\begin{keybox}
\textbf{30-second exam pitch.} Albuquerque, Fos \& Schroth (2022) provide a structural decomposition of the gains from hedge-fund activism into stock picking (pre-existing information advantage) and value creation (engagement improving fundamentals). Using 13D filings 1993--2014 and a model of the activist's accumulation and engagement choice, they estimate that activists capture 20--30\% returns on their stakes, and that 70--85\% of those gains derive from real value creation, with only 15--30\% from information advantage. Value-creation shares are largest for strategic/governance/offensive campaigns and smaller for financial campaigns. The results directly complement Brav-Jiang-Kim (2015) by quantifying the monitoring-and-engagement contribution whose realizations BJK observe at the plant level. Policy: limiting disclosure windows trades off less stock-picking for less value-creation incentive. AFS is the structural benchmark for the 2020s debate on whether activism is a net-positive governance institution.
\end{keybox}

\end{document}
